% Part: lambda-calculus
% Chapter: syntax
% Section: eta

\documentclass[../../../include/open-logic-section]{subfiles}

\begin{document}

\olfileid{lam}{syn}{eta}
\olsection{تحويل~$\eta$}

توجد علاقة أخرى على حدود~$\lambda$. فقد استعملنا في~\olref[fv]{sec}
المثال~$\lambd[x][(fx)]$، وهو يقبل وسيطًا ويطبق~$f$ عليه. وبعبارة أخرى،
إنه الدالة نفسها التي يمثلها~$f$: فكل من~$\lambd[x][(fx)]N$ و$fN$
يختزل إلى~$fN$. ونستعمل اختزال~$\eta$ (وتمديد~$\eta$) لصوغ هذه الفكرة.

\begin{defn}[تقلّص~$\eta$، $\eredone$]
  \ollabel{defn:beredone}
  \emph{تقلّص~$\eta$} ($\eredone$) هو أصغر علاقة متوافقة على الحدود
  تستوفي الشرط الآتي:
  \[
    \lambd[x][M x] \eredone M \text{ بشرط } x \notin \FV{M}
  \]
\end{defn}

\begin{defn}[اختزال~$\beta\eta$، $\bered$] \ollabel{defn:bered}
  \emph{اختزال~$\beta\eta$} ($\bered$) هو أصغر علاقة انعكاسية ومتعدية
  على الحدود تحتوي~$\bredone$ و$\eredone$، أي إنه يضم قاعدتي الانعكاس
  والتعدي فضلًا على القاعدتين الآتيتين:
  \begin{enumerate}
  \item إذا كان~$M \bredone N$ فإن~$M \bered N$. \ollabel{defn:bered3}
  \item إذا كان~$M \eredone N$ فإن~$M \bered N$. \ollabel{defn:bered4}
  \end{enumerate}
    
\end{defn}

\begin{defn}
  نوسع علاقة التكافؤ~$\equal$ بقاعدة تحويل~$\eta$:
  \[
  \lambd[x][f x] \equal f
  \]
  ونرمز إلى العلاقة الموسعة بـ$\equal[\eta]$.
\end{defn}

للتكافؤ بحسب~$\eta$ أهمية لأنه مرتبط بامتدادية حدود لامبدا:

\begin{defn}[الامتدادية]
  نوسع علاقة التكافؤ~$\equal$ بقاعدة~(\ext):
  \begin{center}
  إذا كان~$Mx \equal Nx$ فإن~$M \equal N$، بشرط~$x \notin \FV{MN}$.
  \end{center}
  ونرمز إلى العلاقة الموسعة بـ$\equal[\ext]$.
\end{defn}

وبعبارة تقريبية، تقرر القاعدة أن حدين، حين ينظر إليهما بوصفهما دالتين،
ينبغي أن يعدّا متساويين إذا كان سلوكهما واحدًا عند الوسيط نفسه.

نبرهن الآن أن قاعدة~$\eta$ تعطي الامتدادية بالضبط، ولا تعطي شيئًا زائدًا.

\begin{thm}
  $M \equal[\ext] N$ إذا وفقط إذا كان~$M \equal[\eta] N$.
\end{thm}

\begin{proof}
  نبرهن أولًا أن~$\equal[\eta]$ مغلقة تحت قاعدة الامتدادية. أي إن قاعدة
  $ext$ لا تضيف شيئًا إلى~$\equal[\eta]$. ومن ثم تحتوي~$\equal[\eta]$
  العلاقة~$\equal[\ext]$، وإذا كان~$M \equal[\ext] N$ فإن~$M
  \equal[\eta] N$.

  ولإثبات أن~$\equal[\eta]$ مغلقة تحت~\ext، لاحظ أنه لأي~$M \equal N$
  مشتق بقاعدة~\ext{}، لدينا~$Mx \equal[\eta] Nx$ مقدمةً. وعندئذ يكون
  $\lambd[x][Mx] \equal[\eta] \lambd[x][Nx]$ بموجب قاعدة من قواعد
  $\equal$، ويعطينا تطبيق~$\eta$ على الطرفين~$M \equal[\eta] N$.

  وبالمثل نبرهن أن قاعدة~$\eta$ محتواة في~$\equal[\ext]$. فلأي
  $\lambd[x][Mx]$ و$M$ يحققان~$x \notin \FV{M}$، لدينا
  $(\lambd[x][Mx])x \equal[\ext] Mx$، فنحصل على
  $\lambd[x][Mx] \equal[\ext] M$ بموجب قاعدة~\ext{}.
\end{proof}

\end{document}
